\section{The HTTP Layer}

\subsection{Middleware}

The middleware is a single function that either refuses a request or passes it on with
the prefix removed.  Its checks, and the order of them, are in the life of a request
(Section~\ref{sec:guard}); this section records the reasons.

\begin{description}
  \item[Host.]  A page on another name that resolves to the loopback address still
        sends its own name in the \code{Host} header, so an exact allowlist of the two
        loopback names with the port defeats DNS rebinding.  The port is fixed at start.
  \item[Origin, Referer, Fetch metadata.]  A browser serializes \code{Origin} as exactly
        scheme, host and port; anything else (a path, credentials, an opaque \code{null}) is
        refused rather than interpreted.  \code{Sec-Fetch-Site} must be \code{same-origin} or
        \code{none} (or absent, for clients that do not send it).  The value
        \code{same-site} is refused because a page on another port of the same host is
        same-site to the browser, so \code{SameSite=Strict} does not stop it.
  \item[Read-only.]  Only \code{GET} and \code{HEAD} reach a handler; the endpoint table
        contains no other method.
\end{description}

\subsection{The session and the prefix}

The launch token and the session value are 256-bit random values, compared in constant
time.  The token is exchanged once, atomically, for the session cookie, at the bare root
(\code{/?token=}); the response redirects to the prefixed root so the token leaves the
address bar and the history.  The launch URL does not name the prefix: it is handed to the
program that opens the browser, whose command line other users can read, and a user who
knew the prefix could have the user's browser send the cookie to a server of theirs on
another port (finding~47).

A cookie is scoped by host and by path, never by port.  A plain \code{Path=/} cookie would be
sent to every other web server on \code{127.0.0.1} that the same browser visits, and a server
that received it could replay it.  The design therefore serves everything under a random
128-bit path prefix chosen at start, sets the cookie's path to it, and refuses (403) any URL not
under it.  The browser then sends the cookie only for URLs under the prefix, which no other
server knows.  The prefix is a scope rather than a credential, since the session cookie is
still required; but it is kept from other users, because knowing it is what would let a
local server receive the cookie.  The
front end forms every URL relative to its own location so that this needs no configuration.

\subsection{The launch comes from the user's own connection (Linux)}
\label{sec:launchowner}

The launch URL is in the command line of the program that opens the browser (and on
Linux of the browser itself, until it first uses the URL), where other users can
read it.  On Linux, \code{exchange} therefore accepts a correct token only from a
connection whose client end belongs to \texttt{fsb}'s own user (\fq{SEC-16}).  The
order is the point: (1) a wrong token is refused with no lookup; (2) a token already
used is refused the same way, with no lookup and no warning, so replaying it costs
nothing and fills no log; (3) only then is the owner looked up, and a refusal leaves
the token unused, so another user who reads the URL and uses it first cannot keep
the user's own browser out; (4) the token is consumed atomically.  Refusals are
reported on standard error (\code{Config.Warn}), at most three times and once more
to say the rest are not, with the reason but never the token.

The lookup (\code{internal/connowner}) takes the server's and the client's address
from the request (\code{http.LocalAddrContextKey} and \code{RemoteAddr}, IPv4-mapped
addresses unmapped) and finds the row of \code{/proc/net/tcp} or \code{tcp6} whose
own address is the client's and whose peer is the server's: the client's end.
\texttt{fsb}'s own end lists the same two addresses the other way round, and would
name \texttt{fsb}'s own user, so matching it would let anyone through; the test
compares the matched row's socket inode with the client's.  Rows in ESTABLISHED,
FIN\_WAIT1 or FIN\_WAIT2 count (a client may close its sending side after its
request); TIME\_WAIT rows carry no owner.  The table's uid column is given in
\texttt{fsb}'s own user namespace, so a sandboxed browser (Flatpak, Snap) should show
as the user (not yet tried, Section~\ref{sec:untested}); the overflow uid, which every user the namespace cannot name shares, is
treated as unknown.  An unreadable table or a missing row is unknown too, and
unknown is refused: failing open would let the check be skipped by any condition
that hides a row.  A port forward counts as the user running it.

The same check applies to every request that carries a valid session cookie, so that a
cookie captured by another user, by whatever route, is refused.  It runs only after
the cookie is validated, so a request without the session costs no lookup, and its
result is remembered per connection (\code{httpguard.ConnContext}, installed as the
server's \code{ConnContext}): it is the owner that is remembered, not a verdict, and
only once found, so a lookup that fails is tried again on the next request rather
than refusing the whole connection; a browser's few kept-alive connections cost a few
lookups in all.  Refused launches and refused sessions have separate warning budgets.  macOS has no
comparable table, so there \code{RequireOwner} is given no lookup and both checks are
off.

\subsection{Headers and framing}

Every response carries the headers of \fq{SEC-9}.  Two endpoints relax framing, deliberately
and narrowly: the fixed Markdown page and the PDF may be framed by the interface's own pages
(\code{X-Frame-Options: SAMEORIGIN}, \code{frame-ancestors 'self'}), and the interface itself
may not.

\subsection{Streaming}

Listings and searches are newline-delimited JSON, flushed as batches are ready, because a
folder can hold hundreds of thousands of entries and the first screen must not wait for the
last.  A listing that breaks after it has begun ends with an error line rather than a status
code, since the status was already sent.  Pagination was rejected: the client sorts, and a
sort needs all entries, so streaming keeps the whole listing available and still paints early.

\subsection{Errors}

One function maps every guard error to a status.  Denied, outside-the-roots, nonexistent,
malformed and not-regular all map to the same 404 with the same body; only \code{--debug} adds
the rule name.  Permission errors on paths that passed every rule map to 403; cloud-only files to
409; wrong kinds to 415; anything unexpected to 500 with the detail in the log only.

\subsection{Streams end explicitly}

A listing and a search are streamed as newline-delimited JSON, and a response that simply
stops (a dropped connection, a server error after the header) would otherwise look like a
folder that was just short.  Both therefore end with a line that says how they ended: a
listing with \code{\{"done":true\}}, a search with \code{\{"done":true,\ldots\}} carrying
how many entries were visible, whether a limit stopped it, and whether part of it could
not be searched (\code{incomplete}); an error mid-stream is a line of its own instead.
The client requires the done line and treats a stream that ends without one as broken,
not as complete.

\subsection{Names on the wire}
\label{sec:wirenames}

A Linux name may be any bytes except \code{/} and NUL, so it need not be UTF-8, while
JSON strings are Unicode: Go's encoder replaces each byte that is not part of valid
UTF-8 with U+FFFD.  A file so named was listed under a name that does not exist, and
the client, which builds every request from the names it was given, could never ask
for it again.  \code{wireName} (\code{internal/server/wire.go}) therefore rewrites every
name and path the server puts in JSON (\fq{API-10}): valid UTF-8 is unchanged, and each
stray byte becomes NUL followed by two uppercase hexadecimal digits.  No name contains
NUL, so the result is never ambiguous, and \code{wireName} is one-to-one.  Copies are
rewritten (\code{wireEntries}, \code{wireMatches}, \code{wireMeta}); the guard's own
values are not touched, and the guard never sees an escape.  \code{wireArchive} also
passes archive members through it, but finds nothing to escape: the archive lister has
already replaced invalid UTF-8 in them (Section~\ref{sec:content}), since a member is
only ever displayed, never requested.

The return trip needs no decoder.  The client's \code{queryPath}
(\code{web/src/lib/format.ts}), the one place a path becomes a request, turns each
escape into the byte itself (\code{\%E9}) and percent-encodes everything else as UTF-8,
so the \code{path} parameter decodes, on the server, to the name's real bytes, which
the guard already handled.  Keeping the decoder out of the server keeps the escape out
of the security core: a request can only ever name real bytes.  On the screen,
\code{displayName} shows an escape as \code{\textlangle 0xE9\textrangle}, in the style
of the markers of \fq{NAM-2}; ``Copy path'' (\code{copyablePath}) writes the path in
the shell's \code{\$'\ldots'} quoting.  The address bar keeps the escaped form, which
\code{encodeURIComponent} carries as \code{\%00E9}.  A download's suggested name is
the name as received (\fq{NAM-2}), escape included; what a browser makes of a NUL in a
file name is its own choice, and has not been tested.

\subsection{Timeouts and shutdown}

The server sets a header-read timeout against slow clients and has no write timeout,
because streams may legitimately last.  On interrupt or termination it stops accepting requests and
gives requests in progress three seconds.
